\chapter{第10套试卷}

\begin{center}
\textbf{FPGA与VHDL 基础试卷（十）}

（满分：100分 \hspace{2em} 考试时间：120分钟 \hspace{2em} 难度：中等）
\end{center}

% ============================================
\section*{一、选择题（每题3分，共18分）}
% ============================================

\begin{enumerate}[label=\arabic*.]

\item FPGA上电后需要从外部存储器加载配置数据（比特流）。以下哪一项\textbf{不是}常见的FPGA配置模式？
\begin{enumerate}[label=\Alph*.]
  \item 主串模式（Master Serial）
  \item 从串模式（Slave Serial）
  \item JTAG模式
  \item UART串口通信模式
\end{enumerate}

\item 关于CPLD与FPGA在工艺方面的典型区别，以下描述正确的是？
\begin{enumerate}[label=\Alph*.]
  \item CPLD通常采用Flash/EEPROM工艺（非易失性），FPGA通常采用SRAM工艺（易失性）
  \item CPLD通常采用SRAM工艺（易失性），FPGA通常采用Flash/EEPROM工艺（非易失性）
  \item 两者均采用SRAM工艺
  \item 两者均采用反熔丝工艺
\end{enumerate}

\item 在VHDL中，关于并发语句与顺序语句的描述，\textbf{正确}的是？
\begin{enumerate}[label=\Alph*.]
  \item PROCESS内部的语句是并发执行的
  \item 并发信号赋值语句可以放在PROCESS内部
  \item 顺序语句（如IF、CASE）只能出现在PROCESS或子程序内部
  \item 结构体的BEGIN\ldots END之间可以直接使用IF语句
\end{enumerate}

\item 关于Moore型与Mealy型有限状态机，以下说法正确的是？
\begin{enumerate}[label=\Alph*.]
  \item Moore型状态机的输出仅与当前状态有关
  \item Mealy型状态机的输出仅与当前状态有关
  \item Moore型与Mealy型的输出逻辑完全相同
  \item Moore型状态机的输出与当前输入和状态均有关
\end{enumerate}

\item VHDL中，逻辑运算符具有不同的优先级。表达式\verb|A AND B OR C|的运算结果是？
\begin{enumerate}[label=\Alph*.]
  \item \verb|(A AND B) OR C|（从左到右结合）
  \item \verb|A AND (B OR C)|（OR优先级高于AND）
  \item 取决于综合工具的实现
  \item 该表达式在VHDL中语法错误，必须使用括号
\end{enumerate}

\item 在VHDL中，关于GENERIC关键字的描述，\textbf{错误的是}？
\begin{enumerate}[label=\Alph*.]
  \item GENERIC用于向实体传递静态参数（如位宽、计数值等）
  \item GENERIC在ENTITY声明中定义，位于PORT声明之前
  \item GENERIC的值在仿真运行过程中可以被动态修改
  \item 使用GENERIC可以实现参数化设计，提高代码复用性
\end{enumerate}

\end{enumerate}

\newpage

% ============================================
\section*{二、填空题（每题3分，共12分）}
% ============================================

\begin{enumerate}[label=\arabic*.]

\item VHDL中三种重要对象的声明位置不同：SIGNAL在\underline{\hspace{3cm}}（结构体的声明区）中声明；VARIABLE只能在\underline{\hspace{3cm}}（PROCESS或子程序内部）声明；CONSTANT可在\underline{\hspace{3cm}}中声明。

\item 使用GENERIC实现参数化设计时，在ENTITY中声明GENERIC的语法为：\\[3pt]
\verb|GENERIC ( 参数名 : 数据类型 |\underline{\hspace{2cm}}\verb| 默认值 )；|\\[3pt]
\noindent（填一个赋值符号，由冒号和等号组成）。

\item 在结构体（ARCHITECTURE）中，实例化一个已声明的元件（COMPONENT）使用的语句是\underline{\hspace{4cm}}。

\item VHDL中两种并发生成语句分别是\underline{\hspace{2cm}} GENERATE语句（用于重复实例化）和\underline{\hspace{2cm}} GENERATE语句（用于条件实例化）。

\end{enumerate}

\newpage

% ============================================
\section*{三、程序改错题（共5分）}
% ============================================

下面的VHDL代码意图设计一个带异步复位的4位二进制加法计数器，但代码中存在\textbf{5处错误}。请指出每处错误的位置（行号）并写出正确的修改方案。

\begin{lstlisting}[language=VHDL, numbers=left, caption={带错误的4位计数器}]
LIBRARY ieee;
USE ieee.std_logic_1164.ALL;
                                        -- 缺少算术运算库声明
ENTITY cnt4_err IS
  PORT(
    clk, rst : IN  STD_LOGIC;           -- rst: 低电平异步复位
    q        : OUT STD_LOGIC_VECTOR(3 DOWNTO 0)
  );
END cnt4_err;

ARCHITECTURE behav OF cnt4_err IS
  SIGNAL count : STD_LOGIC_VECTOR(3 DOWNTO 0);
BEGIN
  PROCESS(clk)                          -- 敏感信号表不完整
  BEGIN
    IF rst = '0' THEN                   -- 异步复位逻辑
      count := "0000";                  -- 赋值符号错误
    ELSIF rising_edge(clk) THEN
      count <= count + "0001";          -- 算术运算符使用不当
    END IF;
  END PROCESS;

  q := count;                           -- 并发赋值符号错误
END behav;
\end{lstlisting}

\begin{framed}
\textbf{答题要求：}按以下格式逐条写出每个错误的位置、原因和正确写法（每处1分，共5分）。\\
错误1（第\_\_行）：\_\_\_\_\_\_\_\_\_\_ 改为 \_\_\_\_\_\_\_\_\_\_ \\
错误2（第\_\_行）：\_\_\_\_\_\_\_\_\_\_ 改为 \_\_\_\_\_\_\_\_\_\_ \\
错误3（第\_\_行）：\_\_\_\_\_\_\_\_\_\_ 改为 \_\_\_\_\_\_\_\_\_\_ \\
错误4（第\_\_行）：\_\_\_\_\_\_\_\_\_\_ 改为 \_\_\_\_\_\_\_\_\_\_ \\
错误5（第\_\_行）：\_\_\_\_\_\_\_\_\_\_ 改为 \_\_\_\_\_\_\_\_\_\_
\end{framed}

\newpage

% ============================================
\section*{四、程序阅读分析题（共5分）}
% ============================================

阅读以下VHDL代码，回答问题。

\begin{lstlisting}[language=VHDL, numbers=left, caption={分析用VHDL代码}]
LIBRARY IEEE;
USE IEEE.STD_LOGIC_1164.ALL;
USE IEEE.NUMERIC_STD.ALL;

ENTITY freq_div8 IS
  PORT(
    clk     : IN  STD_LOGIC;
    rst     : IN  STD_LOGIC;
    clk_out : OUT STD_LOGIC
  );
END freq_div8;

ARCHITECTURE rtl OF freq_div8 IS
  SIGNAL cnt : UNSIGNED(2 DOWNTO 0);
BEGIN
  PROCESS(clk, rst)
  BEGIN
    IF rst = '1' THEN
      cnt <= (OTHERS => '0');
    ELSIF rising_edge(clk) THEN
      cnt <= cnt + 1;
    END IF;
  END PROCESS;

  clk_out <= cnt(2);
END rtl;
\end{lstlisting}

\begin{framed}
\textbf{问题：}
\begin{enumerate}[label=(\arabic*)]
  \item 该VHDL代码实现的是什么电路？其功能是什么？（1分）
  \item 若输入时钟\verb|clk|的频率为8\,MHz，输出\verb|clk_out|的频率是多少？请说明计算依据。（1分）
  \item 信号\verb|cnt|的位宽是多少？计数范围是多少？（1分）
  \item 该电路的复位信号是同步复位还是异步复位？请从代码中找到判断依据并说明。（1分）
  \item 若要将分频比从8改为32（即32分频），代码需要做哪些修改？（1分）
\end{enumerate}
\end{framed}

\newpage

% ============================================
\section*{五、综合设计题（共60分）}
% ============================================

% ---------- 设计题1 ----------
\subsection*{设计题1：组合逻辑设计——BCD码至七段数码管译码器（20分）}

设计一个BCD码至七段数码管译码器。输入为4位BCD码\verb|bcd(3 DOWNTO 0)|（仅0--9有效，10--15不会出现），输出为七段数码管驱动信号\verb|seg(6 DOWNTO 0)|，对应段位分配为：seg(6)=a、seg(5)=b、seg(4)=c、seg(3)=d、seg(2)=e、seg(1)=f、seg(0)=g。输出低电平有效（共阳极数码管），即段亮时为'0'，灭时为'1'。

要求：
\begin{enumerate}[label=(\arabic*)]
  \item （5分）列出完整的真值表（输入0--9，输出7位段码）。
  \item （3分）画出顶层模块的端口框图，标注所有输入输出端口名称、方向和位宽。
  \item （10分）用VHDL编写完整的实体和结构体代码，要求使用选择信号赋值语句（\verb|WITH-SELECT-WHEN|）实现组合逻辑功能。当输入为无效BCD码（10--15）时，所有段均熄灭（输出全'1'）。
  \item （2分）代码规范：命名合理、缩进清晰、注释完整。
\end{enumerate}

\begin{framed}
\textbf{提示：}共阳极数码管段码（低有效）参考值——显示“0”：\verb|"1000000"|（a--f亮，g灭），“1”：\verb|"1111001"|，“2”：\verb|"0100100"|，以此类推。seg(6 DOWNTO 0) = \{a, b, c, d, e, f, g\}。
\end{framed}

% ---------- 设计题2 ----------
\subsection*{设计题2：时序逻辑设计——带预置数的8位双向移位寄存器（20分）}

设计一个8位双向移位寄存器，支持左移（向高位移动）和右移（向低位移动），并具有同步预置数功能。

端口定义：
\begin{itemize}
  \item \verb|clk|（输入）：时钟信号，上升沿触发
  \item \verb|rst|（输入）：同步复位，高电平有效，复位时所有输出位清零
  \item \verb|dir|（输入）：移位方向控制，\verb|dir='1'|时左移，\verb|dir='0'|时右移
  \item \verb|load|（输入）：预置数使能，高电平有效
  \item \verb|data_in(7 DOWNTO 0)|（输入）：预置数并行输入
  \item \verb|si|（输入）：串行移位输入（右移时移入最高位MSB，左移时移入最低位LSB）
  \item \verb|q(7 DOWNTO 0)|（输出）：当前移位寄存器的状态
\end{itemize}

要求：
\begin{enumerate}[label=(\arabic*)]
  \item （5分）画出该移位寄存器的功能框图（端口图），并写出其简化功能表（含复位、预置、左移、右移四种操作）。
  \item （10分）用VHDL编写完整的实体和结构体代码，使用PROCESS实现。各控制信号的优先级为：同步复位 > 预置数 > 移位操作。
  \item （5分）若将此8位移位寄存器的每一位用一个1位D触发器（带使能端）元件级联实现，请使用FOR GENERATE语句写出结构描述风格的VHDL代码（要求：先声明一个DFF\_E的COMPONENT，然后使用FOR GENERATE按照移位方向连接各触发器）。
\end{enumerate}

\begin{framed}
\textbf{提示：}右移时，q(6 DOWNTO 0)移入q(7 DOWNTO 1)，q(7)移入si；左移时，q(7 DOWNTO 1)移入q(6 DOWNTO 0)，q(0)移入si。
\end{framed}

% ---------- 设计题3 ----------
\subsection*{设计题3：状态机设计——``1011''序列检测器（Moore型）（20分）}

设计一个Moore型有限状态机，用于检测串行输入序列``1011''。每个时钟上升沿采样一位输入\verb|din|。当连续检测到``1011''序列时，输出\verb|found|输出'1'（持续一个时钟周期），其余时间输出'0'。允许序列重叠检测（即``101011''中应检测出一次``1011''——最后四位）。

示例：若输入序列为\verb|...0 1 0 1 1 0 1 1 0...|，则第5个时钟周期（检测到完整的``1011''）found=1，第8个时钟周期再次found=1。

要求：
\begin{enumerate}[label=(\arabic*)]
  \item （6分）画出完整的Moore型状态转移图，标注每个状态的名称（及含义）、状态编码（自选编码方案）、状态转移条件和各状态的输出值（found）。\\
  建议使用5个状态：S0（初始/未匹配）、S1（已匹配``1''）、S2（已匹配``10''）、S3（已匹配``101''）、S4（已匹配``1011''——检测成功）。
  \item （3分）列出状态转移表，包含：当前状态、输入din、次态、输出found。
  \item （9分）用VHDL编写完整的实体和结构体代码。要求：
  \begin{itemize}
    \item 使用用户自定义枚举类型定义状态（\verb|TYPE state_type IS (S0, S1, S2, S3, S4)|）；
    \item 采用双进程（或三进程）描述风格：一个时序进程描述状态寄存器，一个组合进程描述次态逻辑和输出逻辑；
    \item 时钟上升沿触发，复位\verb|rst|为同步复位（高电平有效），复位后回到S0状态。
  \end{itemize}
  \item （2分）代码规范：状态编码明确、注释清晰、进程划分合理。
\end{enumerate}

\begin{framed}
\textbf{提示：}Moore型状态机的输出仅取决于当前状态，与输入无关。注意处理重叠检测：例如在S3（已匹配``101''）状态下，若输入din=1则进入S4（检测成功）；若输入din=0，则已匹配的后缀``10''可视为新检测序列的前缀，应回到S2状态。
\end{framed}

\vspace{0.5cm}
{\centering
\rule{0.8\textwidth}{0.5pt}
\vspace{0.3cm}

\textbf{—— 试卷结束 ——}

\vspace{0.5cm}}

\endinput
